\section{Abstract}
\label{sec:Abstract}

Die initiale Konfiguration von Advanced Planning and Scheduling (APS) Systemen erfordert umfangreiche, kundenspezifische Anpassungen \citep{stadtler:supply_chain_management2015, pinedo:scheduling2016}. Im Fall von GRACE, dem APS-System der Seitcom GmbH, umfasst dieser Prozess die Definition von sechs Entitätstypen (Materials, Machines, Products, Processes, BOMs, Recipes) mit komplexen referenziellen Abhängigkeiten - bei einer Durchlaufzeit von bis zu acht Wochen und einer Fehlerrate von 10-15 Fehlern pro Konfiguration. Die zentrale Herausforderung liegt in der Fragmentierung des Konfigurationswissens über Azure DevOps Wikis, E-Mail-Korrespondenzen und zehn Jahre implizitem Erfahrungswissen \citep{nonaka:knowledge_creating_company1995, polanyi:tacit_dimension1966}. Bisher existiert kein systematischer Ansatz zur Externalisierung dieses fragmentierten Wissens für GRACE-Konfigurationen. \\
Ziel dieser Arbeit war die Entwicklung einer systematischen Wissensexternalisierungs-Architektur zur Optimierung des GRACE-Konfigurationsprozesses mittels Action Research \citep{lewin:action_research1946, reason:handbook_action_research2001}. \\
Das Ergebnis ist eine Drei-Artefakte-Architektur kombiniert mit einem LLM-gestützten Automatisierungstool.
\\\\
Die Architektur gliedert sich in vier Hauptkomponenten:
\begin{description}
\item[Erhebungsfragebogen] Strukturierte Wissenserfassung mittels 15 Kernfragen nach Prinzipien der Skalenentwicklung \citep{devellis:scale_development2017, mayring:qualitative2015}; reduziert von 96 initialen Fragen durch iterative Validierung.
\item[Entity Mapping Workbook] JSON-Schemas für sechs Entitätstypen basierend auf Entity-Relationship-Modellierung \citep{chen:erm1976}; definiert Feldtypen, Defaults, Constraints und referenzielle Integrität (Product.material.id, BOM.baseQuantity, Process.processingTime.unitId).
\item[Standard Operating Procedures] Drei SOPs (Kundenworkshop, Master Creator-Nutzung, GRACE-Import) mit 6-Punkt-Validierungscheckliste (Vollständigkeit, ID-Konsistenz, Zeiteinheiten, BOM baseQuantity, Product Material, ID-Duplikate).
\item[Master Creator] LLM-gestütztes Tool (qwen2.5-coder:7b) zur automatisierten Entitätsgenerierung mit lokaler Infrastruktur zur Wahrung der Datenhoheit \citep{hummel:data_sovereignty2021, gupta:multilingual2025}; mehrschichtige Validierung nach Scaffolding-Ansatz \citep{sanwal:layered_cot2024}.
\end{description}
Die Evaluation durch ein kontrolliertes Experiment (n=18, standardisiertes Demo-Werk-Szenario) zeigt eine Zeitreduktion von 47 \% (Demo-Werk: 2 Monate $\rightarrow$ 4 Stunden; Kundenkonfiguration: 3,5 Wochen $\rightarrow$ 2 Stunden) und eine Fehlerreduktion von 81 \% (10-15 Fehler $\rightarrow$ 3-5 Fehler pro Konfiguration). Die Vollständigkeit steigt von 60-70 \% (Ist-Analyse analysierter Konfigurationen, gemessen an 6-Punkt-Validierungscheckliste) auf 98,1 \% (Wilson-95-KI: 89,9–99,9) im Experiment. Ergänzend demonstriert die explorative NAS-Crash-Rekonstruktion (n=1, ungeplantes Real-World-Szenario) die Praxistauglichkeit: 80 \% der Demo-Werk-Konfiguration (27 Materials, 8 Machines, 3 Products, 6 Processes, 3 BOMs, 3 Recipes) wurden in 20 \% der ursprünglichen Arbeitszeit rekonstruiert, ermöglicht durch teilweise vorhandene JSON-Konfigurationen als externalisierte Wissensträger. \\\\
Die Ergebnisse validieren den Scaffolding-Ansatz mit mehrschichtiger Validierung und positionieren die Drei-Artefakte-Architektur als effektiven Ansatz für systematische Wissensexternalisierung in APS-Konfigurationsprozessen; die Evaluation ist deskriptiv (kontrolliertes Experiment n=18, explorativ n=1) und belegt Machbarkeit sowie Effizienz, jedoch keine statistische Generalisierbarkeit auf andere APS-Systeme ohne Anpassung der Artefakte.

